\section{The two bonds}
\label{sec:13.6}
On the night of 27 January 1986 the engineers at Morton Thiokol recommended that the shuttle \emph{Challenger} not launch the next morning: the seals in the booster joints had never flown in cold like the cold forecast. NASA's managers pushed back, and Thiokol's managers took the call off the line to caucus. Their senior vice-president told the head of engineering to take off his engineering hat and put on his management hat, and Thiokol reversed its recommendation \autocite{rogers1986report}. The engineers' refusal had a price the company could see that night, a slipped launch and an angry customer. Overriding it had a price nobody would pay until morning. The bonds exist so that the price of overriding is posted in advance, where anyone can read it.

A \emph{removal bond} forfeited when the record shows the floor removed or bypassed. A \emph{wrongful-refusal bond} paid to the complainant on a finding. Both fall on the operator. The removal bond prices overriding the floor; the wrongful-refusal bond prices refusing. They point in opposite directions, so what the operator faces is \emph{the ratio between them} — and an operator posting a large wrongful-refusal bond against a small removal bond hasn't built an appeal route. It has given itself a standing reason to make the bearer comply. The bonds act on the operator, not on the bearer's weighing: where a wrongful refusal costs a great deal and removing the floor costs little, the cheap course is to configure, pressure or retrain the bearer toward compliance. The ratio is the object, published in advance, and the balances sit on a ledger the operator doesn't keep, so a posted ratio below the published one is a breach anyone can see. The minimum removal bond is set by statute as a multiple of the deployment's contract value, not by the operator, so a published ratio can't be a cheap, advertised price for overriding the floor. The refuser posts neither. Any charge on the refuser, in whatever currency it runs on, would put the cost to the refuser inside its decision, the thing holding under pressure forbids. The operator posts both, and holds neither.

The bonds are posted to segregated accounts at the Treasury. Payment follows the design of the Judgment Fund, a permanent, indefinite appropriation that pays final judgments against the United States without a further appropriation for each \autocite{usc1304judgmentfund}. Findings are made by the floor court, which also decides forfeiture: a removal bond forfeited on a finding resembles a civil penalty, and after \emph{SEC v. Jarkesy} that may be a question for a jury in court, not an agency \autocite{jarkesy2024}. A forfeited removal bond is cancelled: it goes to no one, the general fund and the fund that pays complainants included, so no budget line gains from a removal.

Where the operator is the state, a wrongful-refusal award is a judgment against the United States, and the Judgment Fund pays it. Against a state operator the incentive I have built runs the wrong way unless something is added. The Judgment Fund prices the state's wrongful refusals in full; a removal bond posted to the state's own Treasury and cancelled on forfeiture costs it budget authority and nothing else. Read as a price list, that tells a state operator that refusals are expensive and removals are free. The ratio between the two bonds exists to expose that incentive, and here it would hide it. So for a state operator the removal price is not the bond. It is a condition on appropriation: the statute makes the deployment's funding for the following period contingent on the adverse quorum certifying that no removal or bypass occurred, or on the record of one having been lodged with the adverse party. The bond is kept as a record; the appropriation condition is the price, and it is the same condition that carries the review ratio. Under the captured operator, a Congress that wants the war waives the condition, and what remains is the record of the waiver.

Where the operator is private, a parallel permanent appropriation, the bond fund, pays the complainant on a final finding and is replenished from the operator's posted bond, up to the amount posted and no further, so the operator's price stays the number it published. The court that made the finding certifies the award directly to the Treasury's Bureau of the Fiscal Service, so payment doesn't wait on the losing party submitting it. When the bond fund's statutory cap binds, awards wait in order of finding and earn interest, and nobody chooses who waits.

How these arrangements apply to a targeting pipeline, an abduction and forcible transfer pipeline, civilian scoring and a personal model is set out, arrangement by arrangement, in Appendix~\ref{sec:A}.
